% ══════════════════════════════════════════════════════════════
\section{CENA 4 --- O Qubit que Treme: Integral Tripla \hfill \normalfont\small\color{corTempo}12:00 -- 18:00}
% ══════════════════════════════════════════════════════════════

\subsection{4.1 -- Coordenadas esféricas e o jacobiano}
\tempo{12:00 -- 14:00}

\begin{fala}
A nuvem de incerteza vive dentro de uma esfera. Para calcular volumes em 3D,
o sistema de coordenadas mais natural é o \textbf{esférico}: raio~$\rho$,
ângulo polar~$\phi$ e ângulo azimutal~$\theta$.
\end{fala}

\begin{visual}
\faFilm\ Eixos cartesianos aparecem. Depois os eixos esféricos $\rho$, $\phi$, $\theta$
surgem sobrepostos. Uma fatia da esfera fica destacada mostrando o elemento de volume.
\end{visual}

\begin{mathblock}
\textbf{\color{corMath}\faSuperscript\ A integral tripla em coordenadas esféricas:}

\[
  V_{\text{incerteza}}
  = \iiint_{\text{nuvem}}
    \underbrace{\rho^2 \sin\phi}_{\text{jacobiano}}
    \,\mathrm{d}\rho\,\mathrm{d}\phi\,\mathrm{d}\theta
\]

{\small\color{gray}
\textbf{Por que o jacobiano $\rho^2\sin\phi$?} Ao trocar coordenadas cartesianas
$(x,y,z)$ por esféricas $(\rho,\phi,\theta)$, o elemento de volume se transforma.
$\rho^2\sin\phi$ é a correção que garante que a integral continue medindo volume
correto. Sem ele, a conta estaria errada.}
\end{mathblock}

\ferramenta{WolframAlpha -- integral tripla com jacobiano}
\begin{tela}
\faDesktop\ Query:
\begin{lstlisting}[language=bash]
integrate rho^2 * sin(phi), rho=0 to 0.2752, phi=0 to pi, theta=0 to 2pi
\end{lstlisting}
Resposta: $\approx 0{,}0874$ \checkmark
\end{tela}

\subsection{4.2 -- Conectando temperatura e incerteza}
\tempo{14:00 -- 15:00}

\begin{fala}
Na física dura, o calor ataca o tempo de vida do qubit, o famoso \textbf{tempo de relaxação} ($T_1$). A chance de erro real depende de quanto tempo a operação ($\tau$) dura em relação a essa vida útil.
\end{fala}

\begin{tela}
\faDesktop\ \textbf{Física real:} $P_{\text{erro}} = \frac{\tau}{T_1}$
\end{tela}

\begin{fala}
Para modelar como o calor degrada o qubit, usaremos aqui um \textbf{modelo de acoplamento térmico linear}. Na física real, o acoplamento é via matrizes densidade e tempo de relaxação $T_1$, mas este modelo geométrico captura didaticamente a essência da decoerência térmica que estamos estudando hoje.
\end{fala}

\begin{mathblock}
\textbf{\color{corMath}\faSuperscript\ Modelo de acoplamento térmico-quântico:}

\[
  r_T
  \;=\;
  r_{\max}
  \cdot
  \frac{T_{\text{média}} - T_{\text{amb}}}{T_{\text{junc}} - T_{\text{amb}}}
  \;=\;
  0{,}5 \cdot \frac{66{,}28 - 25}{100 - 25}
  \;=\;
  0{,}5 \cdot \frac{41{,}28}{75}
  \qquad
  \boxed{\;r_T \approx 0{,}2752\;}
\]
\end{mathblock}

\begin{fala}
O raio de 0,2752 significa que o calor empurra o estado quântico até 27,5\%
do raio da Esfera de Bloch. Não é desprezível.
\end{fala}

\subsection{4.3 -- Calculando o volume de incerteza}
\tempo{15:00 -- 17:00}

\begin{fala}
Vamos visualizar essa falha como um "volume de incerteza" na Esfera de Bloch. Embora na física rigorosa o ruído térmico seja descrito pela perda de pureza do estado quântico, essa representação geométrica é uma das ferramentas mais poderosas para visualizarmos, através do cálculo, o impacto do calor na fidelidade do qubit.
\end{fala}

\begin{mathblock}
\textbf{\color{corMath}\faSuperscript\ Volume da nuvem (integral tripla):}
\[
  V_{\text{inc}}
  = 2\pi \cdot 2 \cdot \int_0^{r_T}\rho^2\,\mathrm{d}\rho
  = \tfrac{4}{3}\pi\,(0{,}2752)^3
  \qquad
  \boxed{\;V_{\text{inc}} \approx 0{,}0874\;}
\]
{\small \textbf{Nota de rigor:} Esta modelagem geométrica (volume de incerteza) 
corresponde visualmente à perda de pureza do estado quântico ($\rho \to$ estado misto) 
devido ao acoplamento com o reservatório térmico.}
\end{mathblock}

\ferramenta{Visualização 3D: Esfera de Bloch e Nuvem de Incerteza}
\begin{lstlisting}[caption={Código para visualização 3D}]
import numpy as np
import matplotlib.pyplot as plt

fig = plt.figure(figsize=(10, 8))
ax = fig.add_subplot(111, projection='3d')

# 1. Esfera de Bloch (O recipiente sólido e fosco)
u, v = np.mgrid[0:2*np.pi:30j, 0:np.pi:20j]
x = np.cos(u)*np.sin(v)
y = np.sin(u)*np.sin(v)
z = np.cos(v)
ax.plot_surface(x, y, z, color='royalblue', alpha=0.15, shade=True)
ax.plot_wireframe(x, y, z, color="royalblue", alpha=0.2, linewidth=0.5)

# 2. Nuvem de Incerteza (O erro)
rT = 0.2752
n_pontos = 1000
phi = np.random.uniform(0, 2*np.pi, n_pontos)
costheta = np.random.uniform(-1, 1, n_pontos)
u = np.random.uniform(0, 1, n_pontos)
theta = np.arccos(costheta)
r = rT * u**(1/3)
x_n = r * np.sin(theta) * np.cos(phi)
y_n = r * np.sin(theta) * np.sin(phi)
z_n = r * np.cos(theta) + 1 
ax.scatter(x_n, y_n, z_n, color='orange', alpha=0.4, s=2)

ax.set_box_aspect([1,1,1]); ax.view_init(elev=20, azim=45)
plt.show()
\end{lstlisting}

\subsection{4.4 -- A probabilidade de erro}
\tempo{17:00 -- 18:00}

\begin{mathblock}
\textbf{\color{corMath}\faSuperscript\ Probabilidade de bit-flip:}
\[
  P_{\text{erro}} = r_T^3 \approx 2{,}08\%
\]
\end{mathblock}

\begin{fala}
Dois vírgula zero oito por cento. Parece pequeno? Na computação quântica real, a probabilidade de um algoritmo funcionar depois de várias operações despenca exponencialmente.
\end{fala}

\begin{tela}
\faDesktop\ \textbf{Acúmulo de Erros (d operações):}
\[ P_{\text{sucesso}} = (1 - P_{\text{erro}})^d \]
\end{tela}

\begin{fala}
Se você rodar um algoritmo com 1000 portas lógicas ($d = 1000$), a chance de sucesso com esse nosso erro de 2,08\% é de absurdos \textbf{0,00000007\%}. O erro se acumula. O resultado fica completamente inutilizável.
\end{fala}